\section{Introduction}

A neural network encoder trained exclusively on small MLP and CNN checkpoints ---
without ever seeing a transformer during training ---
can predict ViT model accuracy three times better than the state-of-the-art
weight-space self-supervised learning method.
The key is not the symmetry group. It is the coordinate system.

This result emerges from a question at the intersection of two research agendas:
\emph{how do we represent the weights of a neural network in a way that is useful?}
and \emph{can representations learned from one family of architectures transfer to another?}
These questions matter because model zoo analysis --- predicting properties of arbitrary trained
networks without retraining --- is increasingly important for model selection,
architecture search, and continual learning at scale.
Yet existing methods silently assume that the test models come from the same architecture
family as the training models.
When that assumption breaks, they fail.

We find that it breaks in a specific, diagnosable way.
The SANE method~\cite{Schurholt2024Towards} ---
which achieves strong performance ($R^2 = 0.72$) in its intended within-architecture setting ---
represents each neural network as a flattened sequence of weight scalars,
agnostic to network topology.
When presented with ViT checkpoints after training on MLP/CNN checkpoints,
SANE maps all inputs to near-identical latent codes ($\text{std} \approx 0.0014$).
The representation has collapsed: every ViT model looks the same to the encoder,
regardless of its accuracy.
This is not a degradation in quality; it is a fundamental failure of the representation
substrate under cross-architecture distribution shift.

The failure mode reveals the solution.
What flat tokenization discards is the \emph{relational structure} of a neural network ---
the computational graph of neurons connected by scalar weights.
This structure is architecture-agnostic: a ViT attention projection layer,
a CNN convolutional filter, and an MLP linear layer are all, at bottom,
directed graphs of neurons and edges.
If we encode weights \emph{in terms of their graph structure} rather than as flat sequences,
the encoder sees a consistent vocabulary across architectures.
A permutation-equivariant graph neural network trained on this representation
can process ViT weight graphs even though it was trained only on CNN weight graphs,
because the graph schema is the same.

Our key insight:
\textbf{the directed computational graph schema (node=neuron, edge=weight)
provides a universal weight-space coordinate system that enables SSL representations
trained on MLP/CNN model zoos to generalize directly to ViT model zoo property prediction
without any ViT training data.}
The specific symmetry group --- whether to enforce scale equivariance in addition to
permutation equivariance --- turns out to be secondary.
Counter to predictions from supervised same-architecture
settings~\cite{Kalogeropoulos2024Scale},
scale+permutation equivariance (EquiSSL) achieves $R^2 = 0.185$ on held-out ViT models
while permutation-only equivariance (EquiSSL-perm) achieves $R^2 = 0.231$ ---
a reversal of the expected ordering.
We hypothesize that LayerNorm in ViT architectures eliminates the scale degree of freedom
that scale equivariance targets, making scale invariance an incorrect inductive bias for
ViT weight encoding (see Section~\ref{sec:layernorm}).
This LayerNorm gauge-fixing hypothesis is empirically motivated;
independent theoretical treatment, if verified, would provide additional support.

This paper makes the following contributions:

\textbf{(1)} We identify and characterize the latent collapse failure mode of flat tokenization
under cross-architecture SSL.
On 53 real ViT-S/16 ImageNet checkpoints, SANE achieves $R^2 = 0.072$ --- near-chance
for regression --- with latent variance $\text{std} \approx 0.0014$ indicating collapsed
representations.
Note that SANE is not a generally poor method --- it achieves $R^2 = 0.72$ in its intended
within-architecture setting; the $R^2 = 0.072$ reflects the cross-architecture distribution
shift failure mode specifically.

\textbf{(2)} We demonstrate that graph-based SSL with permutation equivariance achieves meaningful
cross-architecture transfer without target-domain training data.
EquiSSL-perm achieves $R^2 = 0.231$ on 53 held-out ViT-S/16 models ---
a ${\sim}220\%$ improvement over SANE, providing initial proof-of-concept evidence for
cross-architecture SSL transfer.

\textbf{(3)} We discover that scale equivariance does not improve over permutation-only equivariance
for ViT cross-architecture transfer, providing empirical support for the LayerNorm
gauge-fixing hypothesis.
EquiSSL achieves $R^2 = 0.185$ (ablation evaluation, Table~\ref{tab:ablation}),
underperforming EquiSSL-perm ($R^2 = 0.327$ in the same ablation evaluation;
$R^2 = 0.231$ in the main experiment, Table~\ref{tab:main}).
The directional ordering (EquiSSL-perm $>$ EquiSSL $>$ SANE) is consistent across both
evaluations, though absolute magnitudes differ due to training duration differences.

\textbf{(4)} We diagnose why graph-based latent interpolation fails as a model editing primitive.
Latent interpolation achieves $\text{acc} = 0.100$ = exactly random chance ($1/10$ for CIFAR-10),
while weight-space averaging achieves $\text{acc} = 0.125$ ($p \approx 0$, Cohen's $d = -1.10$).
Property prediction and model generation are separable capabilities requiring architecturally
distinct decoders.

We organize the paper as follows.
Section~\ref{sec:related} surveys related work.
Section~\ref{sec:method} describes our methodology.
Section~\ref{sec:setup} details our experimental setup.
Section~\ref{sec:results} presents results.
Section~\ref{sec:discussion} discusses surprising findings, limitations, and broader implications.
Section~\ref{sec:conclusion} concludes.
